% ================================================================
%                            前言
% ================================================================
\chapter*{前言}
\addcontentsline{toc}{chapter}{前言}
\markboth{前言}{前言}

这份关于极小曲面的讲义, 其写作动机是希望有一份几乎不设预备知识、却能一直讲
到当前研究课题的处理. 我们也涉及它在其他领域的一些应用, 包括低维拓扑、广义
相对论与材料科学.

极小曲面的历史可以追溯到 Euler 与 Lagrange, 以及变分法的开端. 其中发展出的
许多技巧在几何与偏微分方程中起了关键作用. 例子包括源于极小曲面正则性理论的
单调性与切锥分析、源于 Bernstein 经典工作中基于极值原理的非线性方程估计,
甚至包括 Lebesgue 在其关于极小曲面 Plateau 问题的学位论文中给出的积分定义.

本书唯一需要的预备知识是 Riemann 几何的基本知识, 以及对极值原理的某种熟悉.
在接近极小曲面的诸种途径(复分析、偏微分方程、几何测度论)之中, 我们选择侧重
理论的偏微分方程方面.

在第 1 章中, 我们首先把极小曲面方程作为图上面积泛函的 Euler--Lagrange 方程
推导出来. 随后我们导出极小曲面方程的参数形式(第一变分公式). 第 1 章的重点
是极小曲面的基本性质, 包括面积的单调性公式与 Bernstein 定理. 我们也提到一些
例子. 在第 1 章的倒数第二节中, 我们导出第二变分公式、稳定性不等式, 并定义
极小曲面的 Morse 指标. 在最后一节中, 我们引入多值极小图, 它在后面讨论
\cite{CM3}--\cite{CM7} 的结果时将起主要作用. 我们还会给出一个来自 \cite{CM18}
的螺旋型极小曲面的局部例子(像螺旋面那样), 它可以分解为多值图, 但螺旋的速率
远不是常数.

第 2 章处理 Bernstein 定理的各种推广. 本章开头, 我们对 $\R^n$ 中极小超曲面
$\Sigma$ 的第二基本形式模长的 Laplacian 导出 Simons 不等式. 在随后的几节中,
我们讨论这个不等式的各种应用. 第一个应用是 Choi--Schoen 定理, 它给出总曲率
很小的极小曲面的曲率估计. 利用这个估计, 我们给出 Heinz 关于极小图的曲率估计
的一个简短证明. 接下来, 我们讨论三维流形中稳定极小曲面的先验估计, 包括
Colding--Minicozzi 关于面积与总曲率的估计, 以及 Schoen 的曲率估计. 此后,
我们遵循 Schoen--Simon--Yau 的做法, 把 Simons 不等式与稳定性不等式结合起来,
对稳定极小超曲面的第二基本形式模长平方给出更高的 $L^p$ 界. 这些更高的 $L^p$
界随后又与 Simons 不等式一起用来给出稳定极小超曲面的曲率估计, 并给出第 1 章
所证 Bernstein 定理的一个推广, 它归功于 De Giorgi、Almgren 与 Simons. 我们
引入``几乎稳定''的概念, 它在理解嵌入曲面时起关键作用. 接下来我们回到多值
极小图, 证明 \cite{CM3} 中的一个重要结果: 若多值图有足够多的层, 则层间的
分离度是次线性增长的. 本章最后讨论 Euclidean 空间中的极小锥及其与 Bernstein
定理的关系.

第 3 章开头引入驻定变分流形, 作为经典极小曲面的推广. 接着我们证明
Michael--Simon 的 Sobolev 不等式. 此后, 我们证明上一章所讨论 Bernstein 定理
的一个推广, 它归功于 Colding 与 Minicozzi. 也就是说, 遵循 \cite{CM6}, 我们
将在第 3 章证明: 密度的一个上界实际上给出了极小曲面所在的最小仿射子空间
的一个上界. 我们将从坐标函数(更一般地, 从调和函数)在 Euclidean 空间中任意
余维的 $k$ 维可求长驻定变分流形上的性质导出这个定理. 最后, 在最后一节中,
我们引入另一种弱收敛的概念(称为气泡收敛), 它是为了解释共形不变问题(包括
二维调和映射与 $J$-全纯曲线)中出现的冒泡现象而发展起来的. 我们将证明气泡
收敛蕴含变分流形收敛.

第 4 章从曲面间映射的经典 Plateau 问题的解开始. 在二维情形, 能量与面积之间
有密切的联系, 主要问题在于理解能量有界的映射缺乏紧性, 这称为``冒泡''. 前
三节给出圆盘上映射的 Dirichlet 问题与 Plateau 问题的基本存在性结果, 第四节
讨论分支点. 此后, 我们转向二维球面上调和映射的存在性, 遵循 Sacks 与
Uhlenbeck \cite{SaUh} 的方法: 先极小化一个扰动后的能量泛函, 再令扰动趋于零
取极限.

在第 5 章中, 我们用一个非常一般的论证来在任意球面上寻找极小球面, 其基本想法
可追溯到 1870 年代的 H. A. Schwarz 与 1910 年代的 G. D. Birkhoff. 处理方式
遵循 \cite{CM27} 与 \cite{CM28}, 其中一些存在性结果是新的. Schwarz 与
Birkhoff 的想法都是用极小极大论证来证明变分问题临界点的存在性. 这特别使得
我们能够构造不稳定的极小曲面. 在极小曲面的极小极大构造中, 用曲面族扫出流形,
并记录扫出中每一片曲面的面积; 然后试图提取一个收敛的极大片序列, 使极大片的
面积收敛到所有扫出中极大片面积的下确界.

第 6 章集中于 Plateau 问题经典解的正则性. 在关于唯一延拓与节点集的一般讨论
之后, 我们研究极小曲面在分支点或非横截相交点附近的局部描述. 遵循 Osserman
与 Gulliver, 我们排除 Plateau 问题解的内部分支点. 在本章余下部分, 我们证明
当边界位于平均凸区域的边界中时, Plateau 问题的解是嵌入的. 这最后一个结果
归功于 Meeks 与 Yau.

在第 7 章中, 我们讨论三维流形中的极小曲面理论. 我们首先说明如何把先前的结果
推广到这一情形(特别是单调性、强极值原理, 以及极小曲面的其他一些基本估计).
然后我们证明 Hersch 以及 Choi--Wang 的结果. 接下来, 我们证明 Ricci 曲率为正
的三维流形中嵌入极小曲面的 Choi--Schoen 紧性定理. 这个紧性结果的一个要点是:
由 Choi--Wang 与 Yang--Yau 的结果, 这类极小曲面具有一致的面积界. 此后我们
证明 Schoen--Yau 的正质量定理. 在最后一节中, 我们证明同伦三维球面上 Ricci
流的 Colding--Minicozzi 有限灭绝定理.

最后, 在第 8 章中, 我们介绍 $\R^3$ 中嵌入极小曲面的一些近期结果. 我们从一个
来自 \cite{CM4} 的局部结果开始: 它说明一个嵌入的极小圆盘要么是图, 要么在稍大
一点的尺度上含有一个双螺旋阶梯. 我们还陈述来自 \cite{CM6} 的单侧曲率估计.
这个定理大致断言: 位于平面一侧且靠近该平面的 $\R^3$ 中的嵌入极小圆盘, 是该
平面上的图. 这个估计的新颖之处在于, 它不像本书前面各章讨论的所有经典极小
曲面结果那样需要任何先验界. 这些结果是 \cite{CM3}、\cite{CM4}、\cite{CM5} 与
\cite{CM6} 一系列论文中所得到的嵌入极小曲面结构理论的起点. 此后, 我们描述
单侧曲率估计的一个应用(来自 \cite{CM10}), 即证明广义 Nitsche 猜想(最初由
P. Collin \cite{Co} 证明). 再往后, 我们转向 \cite{CM24} 中关于嵌入曲面的
Calabi--Yau 猜想的解决. 最后, 我们描述 \cite{CM7} 中关于有限亏格嵌入极小
曲面的主要结构定理, 以及若干依赖于 \cite{CM3}--\cite{CM7} 结构理论的近期
唯一性结果.

本书末尾给出问题与习题.

感谢 Matthias Weber 为本书绘制了部分插图, 也感谢 Sergei Gelfand 的坚持与鼓励.

在本书写作期间, 两位作者均受到来自美国国家科学基金会的常规资助与一项 FRG
资助的支持. 文中表达的任何观点、发现、结论或建议均属作者本人, 未必反映美国
国家科学基金会的观点.

\cleardoublepage
